\chapter{Methodik}
\label{cha:methodik}

Um die in Abschnitt~\ref{sec:forschungsfrage} formulierte Forschungsfrage empirisch zu beantworten,
bedarf es einer methodischen Vorgehensweise, die sowohl die wahrgenommene als auch die tatsächlich
gezeigte Reflexion der Studierenden erfasst. Empirische Untersuchungen in der
Software-Engineering-Forschung stützen sich dafür typischerweise auf eine Kombination aus Experimenten,
Fallstudien und Befragungen \cite{stolGuidelinesConductingSoftware2020}. Die vorliegende Arbeit folgt
diesem Grundgedanken, indem sie ein Mini-Experiment im Rahmen einer regulären Lehrveranstaltung mit
einer Befragung sowie einer qualitativen Auswertung der entstandenen Bearbeitungen verbindet.

Dieses Kapitel gliedert sich entsprechend in sieben Abschnitte: Zunächst wird das übergeordnete
Forschungsdesign erläutert (Abschnitt~\ref{sec:forschungsdesign}), gefolgt von der Datenerhebung
(Abschnitt~\ref{sec:datenerhebung}) und der Stichprobe (Abschnitt~\ref{sec:stichprobe}). Anschließend
werden der Aufbau des Mini-Experiments sowie der eingesetzten Fragebögen im Detail dargestellt
(Abschnitt~\ref{sec:fragebogen-experiment}) und das Verfahren der Datenanalyse beschrieben
(Abschnitt~\ref{sec:datenanalyse}), bevor abschließend die Gütekriterien
(Abschnitt~\ref{sec:guetekriterien}) sowie forschungsethische und datenschutzrechtliche Aspekte
(Abschnitt~\ref{sec:forschungsethik}) eingeordnet werden.

\section{Forschungsdesign}
\label{sec:forschungsdesign}
Die vorliegende Arbeit verfolgt einen explorativen Mixed-Methods-Ansatz \cite{kuckartzMixedMethodsMethodologie2014},
um den in Abschnitt~\ref{sec:forschungsfrage} formulierten Vergleich zwischen dem Reflexionsbot und
einem generischen KI-System empirisch zu untersuchen. Qualitative und quantitative Daten werden dabei nicht nacheinander, sondern parallel
erhoben und erst in der Auswertung zusammengeführt (konvergentes Parallel-Design, bei Creswell und
Plano Clark \emph{convergent design})
\cite{kuckartzMixedMethodsMethodologie2014, creswellDesigningConductingMixed2018}: Die quantitativen Fragebogendaten
liefern ein breites, aber oberflächliches Bild der wahrgenommenen Wirkung beider Systeme, während
die qualitative Inhaltsanalyse der tatsächlichen Bearbeitungen zeigt, wie tief die Studierenden sich
inhaltlich mit den Prozessentscheidungen auseinandergesetzt haben. Erst die Kombination beider
Perspektiven erlaubt es, wahrgenommene und tatsächliche Reflexionstiefe zu unterscheiden.
Abbildung~\ref{fig:forschungsdesign} veranschaulicht diesen Ablauf.

\begin{figure}[htbp]
\centering
\begin{tikzpicture}[
	node distance=0.9cm and 1.1cm,
	box/.style={rectangle, draw, rounded corners, align=center, minimum height=0.9cm, text width=3.6cm, font=\small, inner sep=5pt},
	erhebung/.style={box, draw=black!55, fill=black!7},
	gruppeRB/.style={box, draw=codeblue, line width=0.8pt, fill=codeblue!12},
	gruppeAI/.style={box, draw=codegreen, line width=0.8pt, fill=codegreen!12},
	analyse/.style={box, draw=black!55, fill=white},
	zusammen/.style={box, draw=black, line width=1pt, fill=black!12},
	arr/.style={-{Latex[length=2mm]}, draw=black!65},
	arrRB/.style={-{Latex[length=2mm]}, draw=codeblue, line width=0.8pt},
	arrAI/.style={-{Latex[length=2mm]}, draw=codegreen, line width=0.8pt}
	]
	\node[erhebung] (pre) {Prä-Survey (n=22)};
	\node[gruppeRB, below left=of pre] (rb) {Reflexionsbot-Gruppe: Mini-Experiment (16 Bearbeitungen)};
	\node[gruppeAI, below right=of pre] (ai) {Generic-AI-Gruppe: Mini-Experiment (15 Bearbeitungen)};
	\node[erhebung, below=2.5cm of pre] (post) {Post-Survey (n=20) und Bearbeitungen/Chatlogs};
	\node[analyse, below left=of post] (quant) {Quantitative Analyse (deskriptiv)};
	\node[analyse, below right=of post] (qual) {Qualitative Inhaltsanalyse};
	\node[zusammen, below=2.5cm of post] (merge) {Zusammenführung (Triangulation)};

	\draw[arrRB] (pre) -- (rb);
	\draw[arrAI] (pre) -- (ai);
	\draw[arrRB] (rb) -- (post);
	\draw[arrAI] (ai) -- (post);
	\draw[arr] (post) -- (quant);
	\draw[arr] (post) -- (qual);
	\draw[arr] (quant) -- (merge);
	\draw[arr] (qual) -- (merge);
\end{tikzpicture}
\caption{Konvergentes Parallel-Design des Mini-Experiments}
\label{fig:forschungsdesign}
\end{figure}

\section{Datenerhebung}
\label{sec:datenerhebung}
Die Datenerhebung erfolgte in drei aufeinanderfolgenden Phasen, die alle im Rahmen derselben
Übungseinheit der Lehrveranstaltung \glqq Grundlagen des Software Engineering 1\grqq{} stattfanden.

In der ersten Phase füllten alle Studierenden vor Beginn der eigentlichen Aufgabenbearbeitung einen
Prä-Survey aus. Dieser wurde als freiwilliges, pseudonymisiertes Online-Formular (persönlicher,
selbst gewählter Code statt Klarname) bereitgestellt und erfasste
unter anderem die Vorerfahrung mit KI-Systemen sowie das selbsteingeschätzte Reflexionsverhalten (vgl.
Abschnitt~\ref{sec:fragebogen-experiment}).

In der zweiten Phase bearbeiteten die Studierenden die beiden Prozessszenarien mit dem ihrer Gruppe
zugewiesenen KI-System. Dabei entstanden zwei Arten von Rohdaten: zum einen die schriftlichen
Bearbeitungsdokumente, in denen die Studierenden ihre Entscheidungen und Begründungen zu den
Dilemma-Situationen festhielten, zum anderen die Chatverläufe (Interaktionsprotokolle) mit dem
jeweiligen KI-System. Beide Datenquellen bildeten gemeinsam die Grundlage für die spätere qualitative
Inhaltsanalyse (vgl. Abschnitt~\ref{sec:datenanalyse}).

In der dritten Phase füllten die Studierenden im direkten Anschluss an die Bearbeitung einen
Post-Survey aus, der spiegelbildlich zum Prä-Survey aufgebaut war und die Erfahrung mit dem jeweils
genutzten System sowie die wahrgenommene Unterstützung eigenständigen Denkens erfasste.

Ein persönlicher, selbst gewählter Code erlaubte es, die Angaben aus Prä- und Post-Survey einer Person
zuzuordnen, ohne dass Klarnamen erhoben wurden (Pseudonymisierung). Da dieser Code
teilweise auch als Pseudonym für den Upload der Bearbeitungsdokumente verwendet wurde, konnten die
quantitativen Survey-Daten grundsätzlich auch mit den qualitativ kodierten Fällen derselben Person
verknüpft werden (vgl. Abschnitt~\ref{sec:datenanalyse}).

\section{Stichprobe}
\label{sec:stichprobe}
Die Stichprobe besteht aus Studierenden des Studiengangs Software Engineering (Bachelor of Science) an
der Hochschule Heilbronn in Sontheim, die im Rahmen der regulären Übung zur Lehrveranstaltung \glqq
Grundlagen des Software Engineering 1\grqq{} an dem Mini-Experiment teilnahmen. Diese Veranstaltung ist
im ersten Fachsemester des Studiengangs verortet, was die im Folgenden beschriebene Zusammensetzung der
Stichprobe (überwiegend Erstsemester mit noch geringer einschlägiger Vorerfahrung) erklärt. Die
Teilnahme war freiwillig, die Zuteilung zur Reflexionsbot- oder zur Generic-AI-Gruppe erfolgte durch
die Übungsleitung vor Beginn der Bearbeitung.

Vor der Bearbeitung wurden alle Studierenden im Prä-Survey über Zweck und Ablauf des Mini-Experiments
informiert und um ihre freiwillige Einwilligung zur Teilnahme sowie um die Bestätigung der
Volljährigkeit gebeten. Eine Person verneinte die Einwilligungsfrage ausdrücklich, füllte im Anschluss
jedoch dennoch sowohl Prä- als auch Post-Survey vollständig aus. Da für diese Person keine wirksame
Einwilligung vorliegt, wurden ihre Angaben vollständig von der quantitativen Auswertung ausgeschlossen.
Dieselbe Person konnte ohnehin auch für die qualitative Inhaltsanalyse nicht berücksichtigt werden, da
keine eigene Ausgangsantwort zu den Prozessszenarien auffindbar war.

Für die quantitative Auswertung liegen nach diesem Ausschluss n=20 vollständig ausgefüllte
Post-Survey-Datensätze vor (10 Reflexionsbot, 10 Generic AI) sowie n=22 Prä-Survey-Datensätze. Da nicht
jeder Personal Code eindeutig einer Person zuzuordnen war, weil er teilweise doppelt vergeben wurde,
konnten für den Prä-Post-Abgleich n=16 Personen eindeutig verknüpft werden (vgl.
Abschnitt~\ref{sec:datenanalyse}).

Tabelle~\ref{tab:stichprobe} fasst die Zusammensetzung der Stichprobe zusammen,
Abbildung~\ref{fig:stichprobe} stellt die drei kategorialen Merkmale zusätzlich grafisch dar. Die
Stichprobe besteht überwiegend aus Studierenden des ersten Fachsemesters, ergänzt um einzelne
Studierende höherer Semester sowie eine Austauschstudentin bzw. einen Austauschstudenten. Der
überwiegende Teil der Teilnehmenden ist männlich, und der Großteil hatte vor dem Mini-Experiment noch
keine oder nur geringe berufliche Erfahrung im IT-Bereich. Erfahrung mit generativen KI-Chatbots war
dagegen nahezu durchgängig vorhanden. In allen drei erfassten Merkmalen ist die Stichprobe damit
deutlich einseitig zusammengesetzt.

\begin{table}[htbp]
\centering
\caption{Zusammensetzung der Stichprobe (Prä-Survey, n=22)}
\label{tab:stichprobe}
\renewcommand{\arraystretch}{1.3}
\setlength{\tabcolsep}{8pt}
\begin{tabular}{>{\raggedright\arraybackslash}p{4.6cm}
                >{\raggedright\arraybackslash}p{6.0cm}
                >{\raggedleft\arraybackslash}p{1.8cm}}
\toprule
\textbf{Merkmal} & \textbf{Ausprägung} & \textbf{Anzahl} \\
\midrule
Fachsemester        & erstes Fachsemester        & 15 \\
                    & höheres Fachsemester       & 6 \\
                    & Austauschstudierende       & 1 \\
\midrule
Geschlecht          & männlich                   & 19 \\
                    & weiblich                   & 2 \\
                    & keine Angabe               & 1 \\
\midrule
Vorerfahrung mit generativen KI-Chatbots
                    & bereits genutzt            & 20 \\
                    & nicht genutzt              & 2 \\
\midrule
Alter               & 19 bis 29 Jahre (M=22.0)   & \\
Berufserfahrung im IT-Bereich
                    & Median 0 Monate            & \\
\bottomrule
\end{tabular}
\end{table}

\begin{figure}[htbp]
\centering
\begin{tikzpicture}
\begin{axis}[
	xbar stacked,
	width=0.86\textwidth, height=5.2cm,
	xmin=0, xmax=22.4,
	xtick={0,5,10,15,20},
	xlabel={Anzahl Teilnehmende (n=22)},
	xlabel style={font=\small},
	symbolic y coords={KI-Vorerfahrung, Geschlecht, Fachsemester},
	ytick=data,
	y tick label style={font=\small},
	x tick label style={font=\small},
	bar width=20pt,
	axis lines*=left,
	enlarge y limits=0.45,
	every node near coord/.append style={font=\scriptsize, anchor=center},
]
\addplot+[fill=codeblue!75, draw=codeblue,
          nodes near coords, point meta=explicit symbolic,
          every node near coord/.append style={text=white, font=\scriptsize}]
  coordinates {(15,Fachsemester)[erstes Semester: 15]
               (19,Geschlecht)[männlich: 19]
               (20,KI-Vorerfahrung)[bereits genutzt: 20]};
\addplot+[fill=codeblue!32, draw=codeblue,
          nodes near coords, point meta=explicit symbolic,
          every node near coord/.append style={text=black, font=\scriptsize}]
  coordinates {(6,Fachsemester)[höheres Sem.: 6]
               (2,Geschlecht)[2]
               (2,KI-Vorerfahrung)[2]};
\addplot+[fill=codeblue!10, draw=codeblue,
          nodes near coords, point meta=explicit symbolic,
          every node near coord/.append style={text=black, font=\scriptsize}]
  coordinates {(1,Fachsemester)[1]
               (1,Geschlecht)[1]
               (0,KI-Vorerfahrung)[]};
\end{axis}
\end{tikzpicture}
\caption[Zusammensetzung der Stichprobe]{Zusammensetzung der Stichprobe nach Fachsemester,
Geschlecht und Vorerfahrung mit generativen KI-Chatbots (n=22). Die mittleren Segmente stehen für
höhere Fachsemester (6), weibliche Teilnehmende (2) und fehlende KI-Vorerfahrung (2). Die hellen
Segmente stehen für Austauschstudierende (1) und eine fehlende Angabe zum Geschlecht (1).}
\label{fig:stichprobe}
\end{figure}

Für die qualitative Inhaltsanalyse liegen aus den hochgeladenen Bearbeitungsdokumenten und
Chatverläufen 31 auswertbare Analyseeinheiten vor, die sich auf beide Gruppen verteilen (15 Team AI,
16 Reflexionsbot). Ausgeschlossen wurden zwei erkennbare Testeinreichungen (Ordner mit dem Zusatz
\glqq (Test)\grqq{} im Namen) sowie derselbe bereits genannte Fall ohne eigene Ausgangsantwort und ohne
wirksame Einwilligung. Da bei einem erheblichen Teil der Einreichungen die Autorenschaft nicht eindeutig
den Studierenden zuzuordnen ist (vgl. Abschnitt~\ref{sec:datenanalyse}), reduziert sich die für die
Häufigkeitsauswertung nutzbare Stichprobe auf 11 Fälle mit eindeutig studentischer Autorenschaft.

\section{Fragebogen und Mini-Experiment}
\label{sec:fragebogen-experiment}
Das Mini-Experiment wurde im Rahmen der regulären Übung zur Lehrveranstaltung \glqq Grundlagen des
Software Engineering 1\grqq{} durchgeführt. Die Studierenden wurden in zwei Gruppen eingeteilt: die
\emph{Reflexionsbot}-Gruppe (in den Rohdaten und in den Anhangstabellen \emph{Reflection Bot}), die die Aufgaben mit
dem in dieser Arbeit entwickelten Reflexionsbot bearbeitete, und die \emph{Generic-AI}-Gruppe (in
Tabellen und in der Kodierung auch \emph{Team AI}), die einen generischen KI-Chatbot nutzte, der
Anfragen direkt und ohne systematische Rückfragen beantwortet. Als generisches System durften die Studierenden der Vergleichsgruppe frei
zwischen den öffentlich zugänglichen Consumer-Versionen von ChatGPT (OpenAI), Gemini (Google) und
Claude (Anthropic) wählen, wobei der Zugriff über die jeweiligen Weboberflächen auf den eigenen
Endgeräten der Studierenden (Laptop, Tablet oder Smartphone) erfolgte. Anders als beim Reflexionsbot wurde das
Verhalten dieser Systeme nicht über eine Systemnachricht gesteuert. Die einzige inhaltliche Vorgabe
bestand in der Aufgabenanweisung, die eigenen Entscheidungen in das generische System einzugeben und
um eine kritische Analyse beziehungsweise das stärkste Gegenargument zu bitten (\glqq AI
Analysis\grqq{}, vgl. \hyperref[app:szenarien]{Anhang~A.4}). Welche Modellstufe die einzelnen Systeme zum Erhebungszeitpunkt
(24.\,Juni 2026) jeweils bereitstellten, war nicht kontrollierbar und wurde nicht erfasst. Beide
Gruppen erhielten dieselben zwei Prozessszenarien: eine Produktentscheidung in einem regulierten
ML-Projekt (\glqq SmartDose\grqq{}, vgl. Abschnitt~\ref{sec:konzept-reflexionsbot}) sowie
Koordinationsentscheidungen in einem verteilten Team (\glqq LearnLoop\grqq{}). Zu beiden Szenarien
waren jeweils drei Dilemma-Situationen mit Entscheidung, Begründung und Reflexion zu bearbeiten,
gefolgt von einer Einschätzung der Gegenargumentation des jeweiligen KI-Systems.
Tabelle~\ref{tab:bedingungen} stellt die Rahmenbedingungen beider Gruppen gegenüber und macht
sichtbar, welche Variablen kontrolliert waren und welche nicht.

\begin{table}[htb]
\centering
\caption{Gegenüberstellung der beiden Versuchsbedingungen}
\label{tab:bedingungen}
\renewcommand{\arraystretch}{1.3}
\setlength{\tabcolsep}{8pt}
\begin{tabular}{>{\raggedright\arraybackslash}p{0.22\textwidth}
                >{\raggedright\arraybackslash}p{0.32\textwidth}
                >{\raggedright\arraybackslash}p{0.32\textwidth}}
\toprule
\textbf{Merkmal} & \textbf{Reflexionsbot} & \textbf{Generic AI} \\
\midrule
Zugang & Feste Web-Anwendung (Streamlit) &
Frei gewählte Consumer-Oberflächen auf eigenen Endgeräten \\
\midrule
Sprachmodell & \texttt{gpt-4o-mini}, festgelegt &
ChatGPT, Gemini oder Claude; Modellstufe nicht erfasst \\
\midrule
Steuerung des Verhaltens & Feste Systemnachricht mit sokratischen Rückfragen und adaptiver Tiefe &
Keine Systemnachricht, ungelenkte Interaktion \\
\midrule
Direkte Lösungen & Nicht vorgesehen &
Reguläres Systemverhalten \\
\midrule
Vorgabe zur Gegenargumentation & Im laufenden Dialog enthalten &
Eigene Teilaufgabe \glqq AI Analysis\grqq{} \\
\midrule
Protokollierung & Chatexport in der Anwendung enthalten &
Manuelle Sicherung durch die Studierenden \\
\bottomrule
\end{tabular}
\end{table}

Die Gegenüberstellung zeigt zugleich die Grenze des Vergleichs. Kontrolliert waren auf Seiten des
Reflexionsbots Modell, Systemnachricht und Protokollierung, auf Seiten der Vergleichsgruppe keine
dieser Größen. Der Vergleich ist daher als Kontrast zwischen einem gezielt gestalteten System und
einem typischen alltäglichen KI-Zugang zu verstehen und nicht als kontrollierter Modellvergleich
(vgl. Abschnitt~\ref{sec:limitationen}).

Vor der Übung füllten die Studierenden einen Prä-Survey aus, der unter anderem Vorerfahrung mit
KI-Systemen, selbsteingeschätztes Prozessverständnis (in Anlehnung an die Stufen Beschreiben,
Begründen und Abwägen) sowie allgemeines Reflexionsverhalten erfasste. Im Anschluss an die Übung
folgte ein Post-Survey mit spiegelbildlichen Items zur Erfahrung mit dem jeweils zugewiesenen
KI-System, zur wahrgenommenen Unterstützung eigenständigen Denkens und zur eingeschätzten
Kompetenzentwicklung. Beide Fragebögen wurden über ein pseudonymisiertes, freiwilliges Online-Formular
erhoben. Ein persönlicher Code erlaubte die Zuordnung von Prä- und Post-Angaben, ohne dass Klarnamen
erhoben wurden. Abbildung~\ref{fig:ablauf-experiment} fasst den zeitlichen Ablauf der Übungseinheit
mit den jeweils eingesetzten Instrumenten zusammen.

\begin{figure}[htb]
\centering
\begin{tikzpicture}[
	phase/.style={rectangle, draw=black!55, fill=black!8, rounded corners, align=center,
	              text width=3.5cm, minimum height=0.75cm, font=\small\bfseries, inner sep=4pt},
	instr/.style={rectangle, draw=black!40, fill=white, rounded corners, align=left,
	              text width=3.5cm, font=\scriptsize, inner sep=4pt},
	grp/.style={rectangle, rounded corners, align=center, text width=3.3cm,
	            minimum height=0.7cm, font=\scriptsize, inner sep=4pt},
	arr/.style={-{Latex[length=2mm]}, draw=black!65, line width=0.7pt}
	]
% Drei Phasen nebeneinander
\node[phase] (p1) at (0,0)   {Phase 1\\[-1pt] Prä-Survey};
\node[phase] (p2) at (5.3,0) {Phase 2\\[-1pt] Intervention};
\node[phase] (p3) at (10.6,0) {Phase 3\\[-1pt] Post-Survey};
\draw[arr] (p1) -- (p2);
\draw[arr] (p2) -- (p3);

% Instrumente unter Phase 1
\node[instr, below=0.45cm of p1] (i1) {%
	Demografie und KI-Vorerfahrung\\
	Wissenstest ($t_0$)\\
	Selbsteinschätzung ($t_0$)\\
	Reflexionsverhalten ($t_0$)};
\draw[arr] (p1) -- (i1);

% Gruppen unter Phase 2
\node[grp, draw=codeblue, fill=codeblue!10, below=0.45cm of p2] (g1)
	{\textbf{Reflexionsbot}};
\node[grp, draw=codegreen, fill=codegreen!10, below=0.25cm of g1] (g2)
	{\textbf{Generic AI}};
\node[instr, below=0.3cm of g2] (i2) {%
	Szenario SmartDose\\
	Szenario LearnLoop\\
	je Task 1--3 und, nur Generic AI, AI Analysis};
\draw[arr] (p2) -- (g1);

% Instrumente unter Phase 3
\node[instr, below=0.45cm of p3] (i3) {%
	Wissenstest ($t_1$)\\
	Selbsteinschätzung ($t_1$)\\
	Reflexionsverhalten ($t_1$)\\
	Systembewertung\\
	Kognitive Last};
\draw[arr] (p3) -- (i3);

\end{tikzpicture}
\caption[Ablauf des Mini-Experiments mit Erhebungsinstrumenten]{Ablauf des Mini-Experiments mit den
je Phase eingesetzten Erhebungsinstrumenten. Wissenstest, Selbsteinschätzung und Reflexionsverhalten
wurden zu beiden Zeitpunkten mit identischen Items erhoben und über den persönlichen Code
personenbezogen verknüpft.}
\label{fig:ablauf-experiment}
\end{figure}

Neben der Auswahl einzelner Items war auch die strukturelle Anlage des Fragebogens bewusst gestaltet.
Prä- und Post-Survey wurden weitgehend spiegelbildlich aufgebaut, mit identischen oder nur leicht
angepassten Items zu beiden Zeitpunkten, um Veränderungen pro Person über die Bearbeitung hinweg
nachvollziehen zu können, statt lediglich zwei unabhängige Momentaufnahmen zu vergleichen. Die Items
zum selbsteingeschätzten Prozessverständnis wurden gezielt in die drei Stufen Beschreiben, Begründen
und Abwägen gegliedert, die bewusst an die im Reflexionsstufenmodell nach Hatton und Smith
(Abschnitt~\ref{sec:reflexionsstufenmodelle}) unterschiedenen Reflexionsniveaus angelehnt sind: Die
Beschreiben-Items entsprechen den unteren, rein deskriptiven Stufen, die Begründen- und
Abwägen-Items den höheren, tatsächlich reflexiven Stufen. Diese Parallelität erlaubt es, die im
Survey erhobene Selbsteinschätzung inhaltlich auf dieselbe theoretische Grundlage zu beziehen wie die
qualitative Kodierung der Bearbeitungen, auch wenn beide Erhebungen methodisch unabhängig
voneinander bleiben. Ergänzend wurde ein kurzer objektiver Wissenstest zu Scrum und Prozessmodellen
mit sechs Multiple-Choice-Fragen aufgenommen, um die selbstberichtete Kompetenzeinschätzung um ein
von der Selbstauskunft unabhängiges Maß zu ergänzen. Jede Frage wird mit einem Punkt bewertet (0 bis 6
Punkte). Die Fragen und der vollständige Bewertungsschlüssel sind in \hyperref[app:fragebogen]{Anhang~A.5}, die Ergebnisse in
\hyperref[app:statistik]{Anhang~A.6} dokumentiert.

Die überwiegende Mehrheit der Items (Erfahrung mit KI-Systemen, selbsteingeschätzte Kompetenz,
Reflexionsverhalten, Erwartungen sowie die Bewertung des jeweiligen KI-Systems) wurde auf einer
5-stufigen Likert-Skala erhoben (1 = starke Ablehnung, 5 = starke Zustimmung). Die Mittelwerte in
Kapitel~\ref{cha:ergebnisse} beziehen sich, sofern nicht anders angegeben, auf diese Skala. Eine
Ausnahme bildet die Selbsteinschätzung der kognitiven Belastung im Post-Survey, die auf einer separaten
5-stufigen Skala von \glqq sehr leicht\grqq{} bis \glqq sehr schwierig\grqq{} erfasst wurde. Diese
Skala wird in Abschnitt~\ref{sec:ergebnisse} gesondert berichtet, da sich ihre Richtung von der
übrigen Zustimmungsskala unterscheidet. Der
vollständige Fragebogeninhalt beider Erhebungszeitpunkte ist in \hyperref[app:fragebogen]{Anhang~A.5} dokumentiert.
Tabelle~\ref{tab:operationalisierung} fasst zusammen, welches Konstrukt zu welchem Zeitpunkt mit
welchem Format erhoben wurde.

\begingroup
\small
\setlength{\tabcolsep}{6pt}
\renewcommand{\arraystretch}{1.3}
\begin{longtable}{>{\raggedright\arraybackslash}p{0.19\textwidth}
                  >{\raggedright\arraybackslash}p{0.11\textwidth}
                  >{\raggedright\arraybackslash}p{0.19\textwidth}
                  >{\raggedright\arraybackslash}p{0.31\textwidth}}
\caption{Operationalisierung der quantitativ erhobenen Konstrukte}
\label{tab:operationalisierung} \\
\toprule
\textbf{Konstrukt} & \textbf{Zeitpunkt} & \textbf{Format} & \textbf{Beispiel-Item} \\
\midrule
\endfirsthead
\toprule
\textbf{Konstrukt} & \textbf{Zeitpunkt} & \textbf{Format} & \textbf{Beispiel-Item} \\
\midrule
\endhead
\bottomrule
\endfoot
KI-Vorerfahrung & Prä & 4 Items, 5-stufig Likert; 1 Item Häufigkeit &
\glqq I critically evaluate answers from AI systems before adopting them.\grqq{} \\
\midrule
Domänenwissen & Prä und Post & 6 Multiple-Choice-Items, 0--6 Punkte &
\glqq What is the purpose of a Product Backlog in Scrum?\grqq{} \\
\midrule
Prozessverständnis, Stufe 1 (Beschreiben) & Prä und Post & 3 Items, 5-stufig Likert &
\glqq I can explain what a process model in software development is.\grqq{} \\
\midrule
Prozessverständnis, Stufe 2 (Begründen) & Prä und Post & 3 Items, 5-stufig Likert &
\glqq I can explain why a particular process model is suitable for a specific project.\grqq{} \\
\midrule
Prozessverständnis, Stufe 3 (Abwägen) & Prä und Post & 2 Items, 5-stufig Likert &
\glqq I can weigh which process model best fits a specific scenario and explain why another does
not.\grqq{} \\
\midrule
Reflexionsverhalten & Prä und Post & 5 Items, 5-stufig Likert &
\glqq I compared multiple options before making a final decision.\grqq{} \\
\midrule
Bewertung des Systems & Post & 8 Items, 5-stufig Likert &
\glqq The AI system prompted more independent thinking than I expected.\grqq{} \\
\midrule
Kognitive Last & Post & 6 Items, 5-stufig (1 = very easy bis 5 = very difficult) &
\glqq How would you rate the overall cognitive effort required by the exercise?\grqq{} \\
\end{longtable}
\endgroup

Die Übersicht macht die Parallelisierung sichtbar. Vier der acht Konstrukte wurden zu beiden
Zeitpunkten mit identischen Items erhoben, sodass sich Veränderungen personenbezogen berechnen
lassen. Die Bewertung des Systems und die kognitive Last sind naturgemäß nur nach der Bearbeitung
erhebbar, die KI-Vorerfahrung nur davor.

\newpage
\section{Datenanalyse}
\label{sec:datenanalyse}
Die gesammelten Daten werden sowohl qualitativ als auch quantitativ analysiert und anschließend
zusammengeführt.

Die quantitativen Fragebogendaten (Prä- und Post-Survey) werden deskriptiv ausgewertet, insbesondere
im Gruppenvergleich zwischen den Studierenden, die mit dem Reflexionsbot arbeiteten, und jenen, die
das generische KI-System nutzten. Aufgrund der kleinen Stichprobengröße wird bewusst auf
inferenzstatistische Tests und Verallgemeinerungsansprüche verzichtet. Die Mittelwertvergleiche dienen
der deskriptiven Einordnung und der Triangulation mit den qualitativen Befunden.

Die schriftlichen Bearbeitungen der beiden Prozessszenarien sowie die Chatverläufe mit dem jeweiligen
KI-System werden mittels strukturierender qualitativer Inhaltsanalyse nach Mayring
\cite{mayringQualitativeInhaltsanalyse2019} ausgewertet. Dabei kommt das in
Abschnitt~\ref{sec:reflexionsstufenmodelle} eingeführte deduktive Kategoriensystem nach Hatton und
Smith \cite{hattonSmithReflectionTeacherEducation1995} zum Einsatz, ergänzt um eine Zusatzstufe zur
Erfassung fehlender Eigenreflexion (Vermeidung, Aufgabenverweigerung). Die Herkunft des Textes wird
davon getrennt über die Dimension Autorenschaft erfasst (vollständiges Kategoriensystem mit
Kodierregeln und Autorenschafts-Indikatoren siehe \hyperref[app:kategoriensystem]{Anhang~A.1}). Als Analyseeinheit dient jeweils die
schriftliche Reflexion einer Teilnehmerin bzw. eines Teilnehmers zu einem der beiden Szenarien
(Aufgabenteil \glqq Reflect\grqq{} sowie die anschließende Einschätzung der KI-Gegenargumentation).

Da ein Teil der eingereichten Antworten erkennbar nicht von den Studierenden selbst, sondern durch
Kopieren der Systemantwort entstanden ist, wird zusätzlich zur Reflexionsstufe die Autorenschaft des
Textes (studentisch, unklar oder KI-generiert) erfasst. Andernfalls würde die Analyse die
Reflexionsfähigkeit des KI-Systems statt der Studierenden abbilden. Die Kodierung erfolgte durch eine
einzelne Person anhand der Rohdaten (Chatlogs, eingereichte Dokumente, Formulierungsvergleiche
zwischen Personen). Eine Zweitkodierung zur Absicherung der Intercoder-Reliabilität wurde im Rahmen
dieser Arbeit nicht durchgeführt und ist als Limitation in Kapitel~\ref{cha: diskussion} zu benennen.

Der Kodierdurchgang verlief in drei Schritten. Zunächst wurde das aus der Literatur abgeleitete
Kategoriensystem an einer Teilmenge der Bearbeitungen probeweise angewendet, um zu prüfen, ob sich die
Stufendefinitionen am vorliegenden Material überhaupt trennscharf anwenden lassen. Dabei zeigte sich,
dass die vier Originalstufen den häufigsten Fall dieser Stichprobe nicht abbilden, nämlich
Bearbeitungen ohne jede eigene inhaltliche Aussage. Das Kategoriensystem wurde daraufhin um die
Zusatzstufe~0 ergänzt und die Stufen~3 und~4 für den Kontext des Software-Prozessmanagements
ausformuliert (vgl. \hyperref[app:kategoriensystem]{Anhang~A.1}). Im zweiten Schritt wurde das
gesamte Material mit dem angepassten System kodiert, im dritten wurden alle Fälle erneut
durchgegangen, um die zwischenzeitlich präzisierten Regeln einheitlich anzuwenden.

Für die Zuordnung galten dabei zwei Entscheidungsregeln. Enthielt eine Bearbeitung Indikatoren
mehrerer Stufen, wurde die höchste Stufe kodiert, die durch eine ausformulierte Textstelle belegt
ist, nicht die im Text überwiegende. Blieben nach Sichtung aller verfügbaren Rohdaten Zweifel an der
Autorenschaft, wurde nicht \glqq KI-generiert\grqq{}, sondern die Zwischenkategorie \glqq
unklar\grqq{} vergeben. Beide Regeln sind konservativ gewählt: Die erste verhindert, dass einzelne
schwache Passagen eine erkennbare Abwägung entwerten, die zweite, dass ein Verdacht als Nachweis
gewertet wird. Wie stark die Kategorie \glqq unklar\grqq{} die zentrale Auslagerungskennzahl prägt,
wird in Abschnitt~\ref{sec:limitationen} eigens diskutiert.

Abbildung~\ref{fig:kodierbaum} stellt den zweistufigen Ablauf dar. Entscheidend ist die Reihenfolge:
Die Autorenschaft wird \emph{vor} der Reflexionsstufe geprüft, damit nicht die Reflexionsfähigkeit des
KI-Systems anstelle der studentischen Leistung kodiert wird.

\begin{figure}[htb]
\centering
\begin{tikzpicture}[
	start/.style={rectangle, draw=black!55, fill=black!8, rounded corners, align=center,
	              text width=4.6cm, minimum height=0.7cm, font=\small, inner sep=4pt},
	pruef/.style={diamond, aspect=2.4, draw=black!60, fill=white, align=center,
	              font=\scriptsize, inner sep=1pt},
	res/.style={rectangle, draw=black!50, rounded corners, align=center, text width=2.5cm,
	            minimum height=0.6cm, font=\scriptsize, inner sep=3pt},
	stufe/.style={rectangle, draw=codeblue, fill=codeblue!10, rounded corners, align=center,
	              text width=5.2cm, minimum height=0.7cm, font=\scriptsize, inner sep=4pt},
	arr/.style={-{Latex[length=1.8mm]}, draw=black!65},
	lbl/.style={font=\scriptsize, inner sep=1.5pt, fill=white}
	]
\node[start] (a) at (0,0) {Schriftliche Analyseeinheit};
\node[font=\scriptsize\itshape, anchor=west] at (-6.6,-0.9) {Stufe 1: Autorenschaft};
\node[pruef] (p) at (0,-1.7)
	{KI-Merkmale vorhanden?\\ Chatlog, Prompt-Export,\\ Textgleichheit};
\draw[arr] (a) -- (p);

\node[res, fill=black!5] (ki) at (-4.4,-3.3) {\textbf{KI-generiert}};
\node[res, fill=black!5] (un) at (0,-3.3)    {\textbf{unklar}};
\node[res, draw=codegreen, fill=codegreen!12] (st) at (4.4,-3.3) {\textbf{studentisch}};
\draw[arr] (p.west) -- ++(-1.6,0) |- (ki.north) node[lbl, pos=0.72] {eindeutig};
\draw[arr] (p.south) -- (un.north) node[lbl, pos=0.5] {Zweifel};
\draw[arr] (p.east) -- ++(1.6,0) |- (st.north) node[lbl, pos=0.72] {keine};

\node[font=\scriptsize, align=center, text width=4.2cm] (aus) at (-2.2,-4.6)
	{fließt nur in die Autorenschafts\-kennzahl ein, keine Stufenkodierung};
\draw[arr, dashed, draw=black!45] (ki.south) -- (aus.north west);
\draw[arr, dashed, draw=black!45] (un.south) -- (aus.north east);

\node[font=\scriptsize\itshape, anchor=west] at (-6.6,-5.9) {Stufe 2: Reflexionstiefe};
\node[stufe] (s) at (4.4,-5.9)
	{Kodierung nach Hatton und Smith,\\ Stufen 0 bis 4};
\draw[arr] (st) -- (s);
\end{tikzpicture}
\caption[Zweistufiger Kodierprozess]{Zweistufiger Kodierprozess je Analyseeinheit. Nur Bearbeitungen
mit studentischer Autorenschaft werden anschließend nach Reflexionsstufen kodiert. Damit wird die
Textqualität von der studentischen Eigenleistung entkoppelt.}
\label{fig:kodierbaum}
\end{figure}

Die Analyseeinheit wurde bewusst auf den Aufgabenteil \glqq Reflect\grqq{} samt anschließender
Einschätzung der KI-Gegenargumentation begrenzt und nicht auf die gesamte Bearbeitung ausgeweitet.
Nur in diesem Teil fordert die Aufgabenstellung beide Gruppen gleichermaßen zu einer Begründung der
eigenen Entscheidung auf, sodass die kodierten Textstellen über beide Bedingungen hinweg vergleichbar
sind. Die Tasks~1 und~2 verlangen dagegen primär Analyse und Entscheidung, sodass dort auftretende
Unterschiede eher die Aufgabenstellung als die Reflexionstiefe abbilden würden.

Die Zusammenführung beider Datenstränge (Triangulation) erfolgt fallbezogen über den persönlichen
Code in Form einer Gegenüberstellung (Joint Display): Für Personen, bei denen sowohl eine
Autorenschafts- und Reflexionsstufen-Kodierung als auch Survey-Antworten vorliegen, werden das
dokumentierte Bearbeitungsverhalten und die Selbstauskunft Fall für Fall nebeneinandergestellt. Ein
Befund gilt dabei als \emph{konvergent}, wenn beide Perspektiven in dieselbe Richtung weisen (etwa
eine zurückhaltende Selbsteinschätzung bei zugleich wenig reflektierter Bearbeitung), und als
\emph{divergent}, wenn sie sich widersprechen (etwa eine hohe Nützlichkeitsbewertung bei nachweislich
kopierter Bearbeitung). Insbesondere die divergenten Fälle sind für die Fragestellung aufschlussreich
und werden in Abschnitt~\ref{sec:selbstauskunft-verhalten} vertieft.
Abbildung~\ref{fig:jointdisplay} zeigt das Schema dieser Zusammenführung.

\begin{figure}[htb]
\centering
\begin{tikzpicture}[
	strang/.style={rectangle, draw=black!50, rounded corners, align=left, text width=4.3cm,
	               minimum height=1.9cm, font=\scriptsize, inner sep=5pt},
	code/.style={rectangle, draw=black!60, fill=black!8, rounded corners, align=center,
	             text width=3.2cm, minimum height=0.7cm, font=\small\bfseries, inner sep=4pt},
	fall/.style={rectangle, rounded corners, align=left, text width=4.6cm,
	             font=\scriptsize, inner sep=5pt, minimum height=1.35cm},
	arr/.style={-{Latex[length=1.8mm]}, draw=black!65}
	]
\node[strang, draw=codegreen, fill=codegreen!8] (q) at (-3.4,0)
	{\textbf{Quantitativer Strang}\\[2pt]
	 Wissenstest, Prä und Post\\
	 $\Delta$ Selbsteinschätzung\\
	 Systembewertung};
\node[strang, draw=codeblue, fill=codeblue!8] (l) at (3.4,0)
	{\textbf{Qualitativer Strang}\\[2pt]
	 Autorenschaft\\
	 Reflexionsstufe 0 bis 4\\
	 Chatverlauf};

\node[code] (c) at (0,-2.3) {Persönlicher Code};
\draw[arr] (q.south) -- (c.north west);
\draw[arr] (l.south) -- (c.north east);

\node[fall, draw=black!45, fill=black!5] (kon) at (-3.4,-4.3)
	{\textbf{konvergent}\\
	 Beide Perspektiven weisen in dieselbe Richtung, etwa geringe Reflexion bei
	 zurückhaltender Selbsteinschätzung};
\node[fall, draw=black!45, fill=black!5] (div) at (3.4,-4.3)
	{\textbf{divergent}\\
	 Die Perspektiven widersprechen sich, etwa hohe Nützlichkeitsbewertung bei
	 nachweislich kopierter Bearbeitung};
\draw[arr] (c.south) -- (kon.north);
\draw[arr] (c.south) -- (div.north);
\end{tikzpicture}
\caption[Schema der Triangulation im Joint Display]{Schema der Triangulation. Beide Datenstränge
werden über den persönlichen Code fallbezogen verknüpft und in konvergente und divergente Befunde
eingeteilt. Die divergenten Fälle werden in Abschnitt~\ref{sec:selbstauskunft-verhalten}
ausgewertet.}
\label{fig:jointdisplay}
\end{figure}

\section{Gütekriterien}
\label{sec:guetekriterien}
Um die Qualität sowohl des quantitativen als auch des qualitativen Vorgehens einzuordnen, werden im
Folgenden Objektivität, Reliabilität und Validität als die klassischen Gütekriterien empirischer
Forschung \cite{doeringForschungsmethodenEvaluationSozial2023} für beide Teile der Untersuchung
betrachtet.

Die \textbf{Objektivität} der qualitativen Inhaltsanalyse wird durch das deduktive, vorab
festgelegte Kategoriensystem mit expliziten Kodierregeln unterstützt \cite{mayringQualitativeInhaltsanalyse2019}
(vgl. Abschnitt~\ref{sec:reflexionsstufenmodelle} sowie \hyperref[app:kategoriensystem]{Anhang~A.1}): Jede Reflexionsstufe ist über
konkrete, am Text prüfbare Merkmale definiert, statt der kodierenden Person einen weiten
Interpretationsspielraum zu lassen. Für den quantitativen Teil trägt die Standardisierung der
Fragebogenitems (identischer Wortlaut und identische Skala für beide Gruppen, vgl. \hyperref[app:fragebogen]{Anhang~A.5}) zur
Durchführungs- und Auswertungsobjektivität bei.

Für die \textbf{Reliabilität} der Kodierung wurde jeder Fall nicht isoliert, sondern anhand mehrerer
zusammengeführter Rohdatenquellen eingeschätzt: Chatverlauf, eingereichtes Bearbeitungsdokument und,
wo einschlägig, ein Abgleich der Formulierungen mit anderen Personen, um beispielsweise wortgleiche
oder erkennbar kopierte Passagen zuverlässig zu identifizieren (vgl.
Abschnitt~\ref{sec:illustration-kategoriensystem}). Diese Zusammenführung mehrerer Rohdatenquellen
stützt allerdings primär die Validität der Einordnung und ersetzt keine unabhängige Zweitkodierung:
Eine Absicherung der Kodierreliabilität durch eine zweite Kodiererin oder einen zweiten Kodierer
(Intercoder-Reliabilität) wurde im Rahmen dieser Arbeit nicht durchgeführt und ist als Einschränkung
der Reliabilität zu werten (vgl. Abschnitt~\ref{sec:limitationen}). Für den quantitativen Teil erhöht die
pseudonyme, selbst administrierte Erhebung über ein Online-Formular die Wahrscheinlichkeit konsistenter, von der
Anwesenheit einer Untersuchungsperson unbeeinflusster Antworten \cite{doeringForschungsmethodenEvaluationSozial2023}.

Die \textbf{Validität} der Untersuchung stützt sich zum einen auf die inhaltliche Gestaltung der
Prozessszenarien, die bewusst an realen regulatorischen Rahmenbedingungen (MDR, FERPA) ausgerichtet
wurden, um eine für Software-Prozessmanagement realistische Entscheidungssituation abzubilden (vgl.
Abschnitt~\ref{sec:konzeption-szenarien}). Zum anderen erhöht die Triangulation von
Survey- und Inhaltsanalysedaten im gewählten Mixed-Methods-Design
(Abschnitt~\ref{sec:mixed-methods-grundlagen}) die Validität der Gesamtinterpretation, da sich
wahrgenommene und tatsächlich gezeigte Reflexion gegenseitig prüfen lassen, statt sich allein auf
eine der beiden Perspektiven zu verlassen. Grenzen des gewählten Vorgehens werden in
Abschnitt~\ref{sec:limitationen} kritisch reflektiert.

\section{Forschungsethik und Datenschutz}
\label{sec:forschungsethik}
Die Teilnahme am Mini-Experiment war für alle Studierenden freiwillig und unabhängig von der
Bewertung der Übung. Vor Beginn der Bearbeitung wurden die Studierenden im Prä-Survey über Zweck und
Ablauf der Untersuchung informiert und um ihre ausdrückliche Einwilligung zur Teilnahme sowie um die
Bestätigung der Volljährigkeit gebeten (vgl. Abschnitt~\ref{sec:stichprobe}). Eine Person, die die
Einwilligung ausdrücklich verneinte, wurde vollständig von der Auswertung ausgeschlossen.

Personenbezogene Klarnamen wurden zu keinem Zeitpunkt erhoben. Stattdessen wählten die Studierenden
einen selbst gewählten, nicht auf ihre Identität rückführbaren Personal Code, über den Prä- und
Post-Survey sowie, soweit als Pseudonym auch für den Dokumenten-Upload verwendet, die qualitativ
kodierten Fälle miteinander verknüpft werden konnten (vgl. Abschnitt~\ref{sec:datenerhebung}). Die
erhobenen Daten wurden ausschließlich zum Zweck dieser Bachelorarbeit verarbeitet und nicht an
Dritte weitergegeben, im Einklang mit den Grundsätzen der Datenminimierung und Zweckbindung
(Art.~5 Abs.~1 lit.~b und c DSGVO) \cite{eu2016dsgvo}.
